\textbf{Implication for evaluation practice.} If the argument of this
paper is correct, symbolic and scripted simulation is necessary but not
sufficient for validating an LLM/VLM robot planner: it can validate
reasoning strategy, task decomposition, and high-level correctness, but by
construction cannot exercise the geometry-grounding failure class of
\S\ref{sec:taxonomy}, because it never asks the model to compute a
quantity a physical sensor could contradict. We suggest that a
geometry-grounding audit -- identifying every numeric quantity a planning
prompt asks the model to compute itself, and asking whether a sensor or a
paired action already determines that quantity more reliably -- belongs
alongside task-complexity benchmarking as a standard step before hardware
deployment, independent of which reasoning strategy or foundation model is
used.

\textbf{What the missing ablation would add.} \S\ref{sec:validation} is
explicit that our deterministic grounding layer is validated
retrospectively, against replayed real-hardware failures and a unit
regression suite, rather than through a live controlled A/B comparison
on hardware. A live ablation -- re-running the task conditions that
originally surfaced these bugs with the grounding layer enabled versus
disabled -- would let us report a quantitative safety/success delta
attributable to the fix specifically, rather than inferring its
necessity from the failures it was designed to close. Two scopes are
designed but neither is run at the time of writing: a restricted version
(LLM-only, the two release-geometry fixes specifically, roughly half a
day including trials) and a complete one (both modalities, the full
taxonomy including the safety-critical TCP-offset/grasp-width path,
roughly two days). We flag the restricted version as the most direct
next step, since it targets the mechanism the evaluation of
\S\ref{sec:evaluation} actually implicates.

\textbf{Case 6 remains the largest open item.} Everything else in this
paper is either closed (five of six bugs) or bounded by a controlled
comparison we can eventually run (the ablation above). Case 6 is
neither: it is a trajectory-level, multi-body contact problem, and
\S\ref{sec:zone-spacing}'s static-spacing mitigation, while a genuine
improvement over an unbounded layout, does not touch the dynamic half of
the problem at all. Closing it plausibly needs one of (a) a
motion-planning-level check of the approach trajectory against
already-placed objects' known footprints, not just their final target
positions, or (b) a settle-time verification step after each release,
before the next pick begins, so a displaced neighbor is caught and
corrected rather than compounding into the next release. We have not
implemented or evaluated either.

\textbf{Generalizability.} Our evidence comes from one manipulator (UR5cb),
one gripper (OnRobot RG2), one camera rig, and one reasoning strategy.
Deliberately, it comes from a \emph{single} underlying model
(\code{kimi-k2.6}) used for both the LLM and VLM arm, rather than one
model per modality -- an earlier configuration with modality-specific
models is what originally produced the now-non-replicating
value-mapping-reasoning finding of \S\ref{sec:computation-not-perception},
which in hindsight looks like it was at least partly a model-identity
confound rather than a pure grounding-modality effect. We do not claim
that the specific six bugs of Table~\ref{tab:taxonomy} recur verbatim on
different hardware; we claim that the \emph{structural pattern} --
geometry computed by a model rather than measured or derived from a
paired action, and multiple real bodies sharing space over time -- is a
property of the architecture (asking an LLM/VLM to output numeric
geometry at all, and placing more than one object in a shared zone), not
of this particular robot, and should be audited for on any comparable
deployment.

\textbf{Grasp width remains unverified.} Of the quantities in
\S\ref{sec:grounding-layer}, grasp width is the one our system still
cannot cross-check against depth in every case: a single top-down depth
sample determines an object's height unambiguously, but not its
left-right extent. Our best-effort edge-inset sampling
(\S\ref{sec:grounding-layer}) reduces, but does not eliminate, this gap;
a depth-informed rectification of the camera view before grounding is a
natural next step we have not yet completed.

\textbf{Reliability of operator-supplied labels.} Our semi-automatic
evaluation protocol (\S\ref{sec:evaluation}) depends on a single
operator's real-time judgment of safety and sub-task completion for each
trial. We did not measure inter-rater reliability with a second, blinded
annotator on a subsample of trials, which would strengthen confidence in
the reported TS/TSR numbers independent of the geometry-grounding
argument this paper makes. This is not a hypothetical risk: an early pass
over this cycle's data marked every trial safe regardless of outcome,
caught only by cross-checking the safety label against the free-text
operator note recorded for the same trial, and corrected before the
numbers in \S\ref{sec:evaluation} were computed. We report the incident
as a concrete illustration of exactly the reliability concern this
paragraph raises, not as a suggestion that it recurs undetected
elsewhere in this dataset.
